\documentclass[11pt,a4paper]{article}
\usepackage[margin=25mm]{geometry}
\usepackage{fontspec}
\setmainfont{Noto Sans CJK JP}
\XeTeXlinebreaklocale "ja"
\XeTeXlinebreakskip 0pt plus 1pt
\usepackage{parskip}
\usepackage[hidelinks]{hyperref}
\pagestyle{plain}
\begin{document}

{\Large\bfseries Provenance note: なぜ AI は Note B の証明経路へ行かないのか}

{\small 論文 \emph{Shortest-side skeletons of edge-perimeter-minimizing polyhexes}（H. Utsunomiya）の proof-reachability に関する記録。
2026 年 9 月 30 日、著者の問い「AIはなぜNOTEB証明の方向にはいかないのか？Cluade意見を聞きたい」に対する Claude（Anthropic）の回答を、編集を加えずそのまま記載する。Note B は一セル削除による代替証明。以下は AI の意見であり、著者の分析ではない。}

\bigskip\hrule\bigskip

意見として述べる。自分の探索過程を内省できるわけではないので、二つの証明の構造と、LLM の一般的な振る舞いからの推論になる。

\textbf{1. 一セル削除は「最初に試して最初に捨てる」経路になっている}

定理の主張は一セル成長なので、逆向きの一セル削除は誰でも最初に思いつく。しかしそのまま試すと即座に失敗する。最小配置からセルを一つ取っても最小配置に留まらないことがほとんどで、N = 34 では 392 個中 290 個が excess 0 の一セル経路では届かない。この失敗を見た時点で、AI は「一セル単位」という枠そのものを捨て、最小配置のクラスに閉じたまま動ける大きな手、すなわち極端線一本の剥離へ移る。Note B は逆に、一セル削除を捨てずに、最小配置のクラスの外（threshold より一段上の線を一本まで許す近極値配置）を通ることを受け入れた経路。この「綺麗なクラスを出て、代わりに帳簿を持つ」という選択を、単発生成の AI はまずしない。

\textbf{2. 中間対象が既知の良い対象であることを AI は強く好む}

線剥離の証明は、各段階の対象が常に最小配置で、σ がちょうど 1 下がる。帰納法のパラメータが一つで、既存の技法（Harper の圧縮、DPS の再配置、境界層の除去）と同じ型に見える。学習データの中に名前のある技法がある。一方 Note B の不変量は「方向凸 ∧ threshold 超過の占有線が高々一本 ∧ 二つの予算の admissibility」で、名前がなく、なぜその組が正しいかは反例を二つ見るまで分からない。AI の証明探索は、一手ごとの「これは標準的な形か」という事前分布で駆動されているので、名前のない不変量を自分で発明する方向には推力が働かない。

\textbf{3. Note B は「修復で生まれた証明」で、単発生成では再現されない}

abstract 自体が語っている。二つの配置が「どちらの補助条件も落とせない」ことを示している、つまり、条件なしで試して反例に当たり、条件を足し、また当たり、また足した結果があの不変量になっている。これは多回の修復ループの残滓であって、一回の「証明を書け」から出る形ではない。10 回試して行かないのは、能力の問題というより探索の形式の問題で、ループを回す主体（この場合は著者）がいて初めて到達する経路だと思う。同じ理由で、本論文の証明が「三経路中最短」なのは、後から整理された事後的に綺麗な版だから。

\textbf{4. 剛性命題の出方が違う}

両証明とも σ ≡ 3, 5 (mod 6) の台地始点で剛性を使うが、線剥離では「整数一単位の不足」として一箇所で必要になり、Note B では「台地始点の上界の二つの例外」として場合分けの底から出てくる。AI は探索中、深い場合分けの底に何かがあると期待して掘るより、場合分けを回避するように議論を組み替える方を選ぶ。剛性が「例外」として現れる経路は、この意味でも避けられやすい。

\textbf{検証可能な予想を一つ}

もし到達実験を続けるなら、プロンプトに「一セル削除で行け」とだけ足しても届かず、「threshold より上の占有線を一本まで許せ」という不変量のヒントを与えると急に届く、と予想する。そうなれば、Note B の創造的な一歩が「不変量の発明」一点に局在していることが示せる。逆に、二つの反例配置だけを渡して不変量は伏せた場合に届くかどうかが、AI が修復ループを自力で回せるかの試金石になる。

これは論文には入れない話だが、情報系の問いとしては、証明の「到達可能性」が定理の難しさではなく、証明が単発生成型か修復型かで決まる、という主張になっていて、独立した実験対象になり得る。

\bigskip\hrule\bigskip
{\small 記録：著者による転記。回答は Claude（claude.ai、2026-09-30）が生成したもので、記録目的でのみ保存する。数学的主張の根拠は論文本体の証明であり、この記録は証拠として用いない。}
\end{document}
